% ===== Trees (15 problems) =====
\section{Trees {\normalsize\color{muted}$\cdot$ 15 problems}}
\markboth{Trees}{}

\begin{problem}{226}{Invert Binary Tree}{Easy}
\summary{Mirror the tree (swap every left and right child).}
\keyidea{Swap children, recurse on both.}
\complexity{O(n) time, O(h) space}
\begin{code}
class Solution:
    def invertTree(self, root: Optional[TreeNode]) -> Optional[TreeNode]:
        if root:
            root.left, root.right = self.invertTree(root.right), self.invertTree(root.left)
        return root
\end{code}
\end{problem}

\begin{problem}{104}{Maximum Depth of Binary Tree}{Easy}
\summary{Number of nodes on the longest root-to-leaf path.}
\keyidea{depth = 1 + max(depth(left), depth(right)).}
\complexity{O(n) time, O(h) space}
\begin{code}
class Solution:
    def maxDepth(self, root: Optional[TreeNode]) -> int:
        if not root:
            return 0
        return 1 + max(self.maxDepth(root.left), self.maxDepth(root.right))
\end{code}
\end{problem}

\begin{problem}{543}{Diameter of Binary Tree}{Easy}
\summary{Longest path (in edges) between any two nodes.}
\keyidea{DFS returns height; at each node update best with left height + right height.}
\complexity{O(n) time, O(h) space}
\begin{code}
class Solution:
    def diameterOfBinaryTree(self, root: Optional[TreeNode]) -> int:
        best = 0
        def height(node):
            nonlocal best
            if not node:
                return 0
            l, r = height(node.left), height(node.right)
            best = max(best, l + r)
            return 1 + max(l, r)
        height(root)
        return best
\end{code}
\end{problem}

\begin{problem}{110}{Balanced Binary Tree}{Easy}
\summary{Is every node's left/right height difference at most 1?}
\keyidea{Bottom-up DFS returning height, or -1 as an 'unbalanced' sentinel.}
\complexity{O(n) time, O(h) space}
\begin{code}
class Solution:
    def isBalanced(self, root: Optional[TreeNode]) -> bool:
        def h(node):
            if not node:
                return 0
            l, r = h(node.left), h(node.right)
            if l < 0 or r < 0 or abs(l - r) > 1:
                return -1
            return 1 + max(l, r)
        return h(root) >= 0
\end{code}
\end{problem}

\begin{problem}{100}{Same Tree}{Easy}
\summary{Are two trees identical in structure and values?}
\keyidea{Both None -\textgreater{} True; one None or values differ -\textgreater{} False; else recurse on both sides.}
\complexity{O(n) time, O(h) space}
\begin{code}
class Solution:
    def isSameTree(self, p: Optional[TreeNode], q: Optional[TreeNode]) -> bool:
        if not p and not q:
            return True
        if not p or not q or p.val != q.val:
            return False
        return self.isSameTree(p.left, q.left) and self.isSameTree(p.right, q.right)
\end{code}
\end{problem}

\begin{problem}{572}{Subtree of Another Tree}{Easy}
\summary{Is subRoot an exact subtree somewhere inside root?}
\keyidea{At every node of root, test sameTree(node, subRoot).}
\complexity{O(m * n) time, O(h) space}
\begin{code}
class Solution:
    def isSubtree(self, root: Optional[TreeNode], subRoot: Optional[TreeNode]) -> bool:
        def same(a, b):
            if not a and not b:
                return True
            if not a or not b or a.val != b.val:
                return False
            return same(a.left, b.left) and same(a.right, b.right)
        if not subRoot:
            return True
        if not root:
            return False
        return (same(root, subRoot) or self.isSubtree(root.left, subRoot)
                or self.isSubtree(root.right, subRoot))
\end{code}
\end{problem}

\begin{problem}{235}{Lowest Common Ancestor of a Binary Search Tree}{Medium}
\summary{Find the LCA of nodes p and q in a BST.}
\keyidea{Both smaller -\textgreater{} go left; both larger -\textgreater{} go right; otherwise this node is where they split.}
\complexity{O(h) time, O(1) space}
\begin{code}
class Solution:
    def lowestCommonAncestor(self, root: TreeNode, p: TreeNode, q: TreeNode) -> TreeNode:
        cur = root
        while cur:
            if p.val < cur.val and q.val < cur.val:
                cur = cur.left
            elif p.val > cur.val and q.val > cur.val:
                cur = cur.right
            else:
                return cur
\end{code}
\end{problem}

\begin{problem}{102}{Binary Tree Level Order Traversal}{Medium}
\summary{Return node values level by level.}
\keyidea{BFS with a queue; process len(queue) nodes per level.}
\complexity{O(n) time, O(n) space}
\begin{code}
class Solution:
    def levelOrder(self, root: Optional[TreeNode]) -> List[List[int]]:
        res, q = [], deque([root] if root else [])
        while q:
            level = []
            for _ in range(len(q)):
                node = q.popleft()
                level.append(node.val)
                if node.left: q.append(node.left)
                if node.right: q.append(node.right)
            res.append(level)
        return res
\end{code}
\end{problem}

\begin{problem}{199}{Binary Tree Right Side View}{Medium}
\summary{Values visible when looking at the tree from the right, top to bottom.}
\keyidea{BFS by level; take the last node of each level.}
\complexity{O(n) time, O(n) space}
\begin{code}
class Solution:
    def rightSideView(self, root: Optional[TreeNode]) -> List[int]:
        res, q = [], deque([root] if root else [])
        while q:
            for i in range(len(q)):
                node = q.popleft()
                if i == 0:
                    res.append(node.val)       # enqueue right first -> first is rightmost
                if node.right: q.append(node.right)
                if node.left: q.append(node.left)
        return res
\end{code}
\end{problem}

\begin{problem}{1448}{Count Good Nodes in Binary Tree}{Medium}
\summary{A node is good if no node on the path from the root to it is larger. Count good nodes.}
\keyidea{DFS carrying the max value seen on the path.}
\complexity{O(n) time, O(h) space}
\begin{code}
class Solution:
    def goodNodes(self, root: TreeNode) -> int:
        def dfs(node, mx):
            if not node:
                return 0
            good = 1 if node.val >= mx else 0
            mx = max(mx, node.val)
            return good + dfs(node.left, mx) + dfs(node.right, mx)
        return dfs(root, root.val)
\end{code}
\end{problem}

\begin{problem}{98}{Validate Binary Search Tree}{Medium}
\summary{Is the tree a valid BST (strict ordering)?}
\keyidea{DFS passing an open interval (low, high) that every node must sit inside.}
\complexity{O(n) time, O(h) space}
\begin{code}
class Solution:
    def isValidBST(self, root: Optional[TreeNode]) -> bool:
        def ok(node, lo, hi):
            if not node:
                return True
            if not (lo < node.val < hi):
                return False
            return ok(node.left, lo, node.val) and ok(node.right, node.val, hi)
        return ok(root, float("-inf"), float("inf"))
\end{code}
\end{problem}

\begin{problem}{230}{Kth Smallest Element in a BST}{Medium}
\summary{Return the k-th smallest value (1-indexed) in a BST.}
\keyidea{In-order traversal visits values in sorted order; stop at the k-th pop.}
\complexity{O(h + k) time, O(h) space}
\begin{code}
class Solution:
    def kthSmallest(self, root: Optional[TreeNode], k: int) -> int:
        st, cur = [], root
        while st or cur:
            while cur:
                st.append(cur)
                cur = cur.left
            cur = st.pop()
            k -= 1
            if k == 0:
                return cur.val
            cur = cur.right
\end{code}
\end{problem}

\begin{problem}{105}{Construct Binary Tree from Preorder and Inorder Traversal}{Medium}
\summary{Rebuild the tree from its preorder and inorder sequences (unique values).}
\keyidea{preorder[0] is the root; its index in inorder splits left/right subtrees. Use a value-\textgreater{}index map.}
\complexity{O(n) time, O(n) space}
\begin{code}
class Solution:
    def buildTree(self, preorder: List[int], inorder: List[int]) -> Optional[TreeNode]:
        idx = {v: i for i, v in enumerate(inorder)}
        pre = iter(preorder)
        def build(l, r):
            if l > r:
                return None
            val = next(pre)
            node = TreeNode(val)
            m = idx[val]
            node.left = build(l, m - 1)
            node.right = build(m + 1, r)
            return node
        return build(0, len(inorder) - 1)
\end{code}
\end{problem}

\begin{problem}{124}{Binary Tree Maximum Path Sum}{Hard}
\summary{Max sum of any path (any start and end node, must be connected).}
\keyidea{DFS returns the best downward 'arm' (clamped at 0); at each node update answer with val + left + right.}
\complexity{O(n) time, O(h) space}
\begin{code}
class Solution:
    def maxPathSum(self, root: Optional[TreeNode]) -> int:
        best = root.val
        def arm(node):
            nonlocal best
            if not node:
                return 0
            l = max(arm(node.left), 0)
            r = max(arm(node.right), 0)
            best = max(best, node.val + l + r)
            return node.val + max(l, r)
        arm(root)
        return best
\end{code}
\end{problem}

\begin{problem}{297}{Serialize and Deserialize Binary Tree}{Hard}
\summary{Convert a tree to a string and back.}
\keyidea{Preorder with 'N' for nulls; deserialize by consuming tokens in the same preorder.}
\complexity{O(n) time, O(n) space}
\begin{code}
class Codec:
    def serialize(self, root):
        out = []
        def dfs(node):
            if not node:
                out.append("N")
                return
            out.append(str(node.val))
            dfs(node.left)
            dfs(node.right)
        dfs(root)
        return ",".join(out)

    def deserialize(self, data):
        vals = iter(data.split(","))
        def dfs():
            v = next(vals)
            if v == "N":
                return None
            node = TreeNode(int(v))
            node.left = dfs()
            node.right = dfs()
            return node
        return dfs()
\end{code}
\end{problem}
